% docs/manual/user/chapters/02-quickstart.tex
% 第 2 章：Quick Start

\chapter{Quick Start}

本章用 10 分钟带你跑通第一个完整评估：从无到有生成一份配置文件、验证它、启动评估、看输出。后续章节会逐项深入；本章只关心"先把流水线跑起来"。

\section{工作流总览}

\openbench{} 的标准工作流是三个命令：

\begin{enumerate}
  \item \cli{openbench init}\quad 交互式向导生成 \file{openbench.yaml}
  \item \cli{openbench check openbench.yaml}\quad 验证配置 + 检查数据
  \item \cli{openbench run openbench.yaml}\quad 启动评估
\end{enumerate}

\begin{tipBox}
如果你已经有现成的配置文件（来自队友、模板、或前一次评估），可以跳过 \cli{init} 直接进入 \cli{check}。
\end{tipBox}

\section{步骤 1：生成最小配置（\cli{openbench init}）}

启动向导：

\begin{minted}{bash}
openbench init
\end{minted}

向导会按 5 个章节交互式询问，所有问题都有合理默认值，回车即可接受：

\begin{exampleBox}[向导对话示例]
\begin{minted}{text}
OpenBench Configuration Wizard

1. Project Settings
  Project name [my-evaluation]: demo
  Output directory [./output]:
  Start year [2004]: 2010
  End year [2010]: 2014

2. Evaluation Variables
  Available variables (by category):
  Bio:
    [1] Burned_Area
    [2] Gross_Primary_Productivity
    ...
  Water:
    [10] Evapotranspiration
    [11] Total_Runoff
    ...
  Select variables (comma-separated numbers, or 'all') [all]: 10,11

3. Reference Data Sources
  Evapotranspiration - available sources:
    [1] GLEAM_v4.2a_LowRes (grid, Month)
    [2] FLUXCOM-X-BASE_LowRes (grid, Month)
  Select for Evapotranspiration [1]:
  Total_Runoff - available sources:
    [1] GRUN_ENSEMBLE_LowRes (grid, Month)
  Select for Total_Runoff [1]:

4. Simulation Models
  Known models:
    [1] BCC_AVIM (10 variables)
    [2] CLM5 (12 variables)
    [3] CoLM2024 (10 variables)
    ...
  Model name (or number from list, empty to finish): 3
  Data root directory for CoLM2024: /data/colm2024
  Label for this run [CoLM2024]:
  Added: CoLM2024
  Model name (or number from list, empty to finish):

5. Options
  Enable comparison? [Y/n]:
  Enable statistics? [y/N]:

Generated openbench.yaml
  Next: openbench check openbench.yaml
\end{minted}
\end{exampleBox}

\subsection{产出的 YAML}

向导写出的 \file{openbench.yaml} 形如：

\begin{minted}{yaml}
project:
  name: demo
  output_dir: ./output
  years:
  - 2010
  - 2014
evaluation:
  variables:
  - Evapotranspiration
  - Total_Runoff
reference:
  sources:
    Evapotranspiration: GLEAM_v4.2a_LowRes
    Total_Runoff: GRUN_ENSEMBLE_LowRes
simulation:
  CoLM2024:
    model: CoLM2024
    root_dir: /data/colm2024
comparison:
  enabled: true
\end{minted}

这就是 \openbench{} 配置的最小可运行结构：4 个顶层 key（\yamlkey{project} / \yamlkey{evaluation} / \yamlkey{reference} / \yamlkey{simulation}）+ 1 个可选 \yamlkey{comparison}。所有未列出的字段使用 \yamlkey{schema} 默认值（详见用户卷附录 A）。

\begin{noteBox}
向导只产出"最小骨架"。要利用 \openbench{} 的全部能力（自定义指标、统计、远程运行、IGBP/PFT 分组等），需要手动编辑配置文件加字段，详见第 3-7 章。
\end{noteBox}

\section{步骤 2：检查配置（\cli{openbench check}）}

\begin{minted}{bash}
openbench check openbench.yaml
\end{minted}

这一步**不读任何模型/参考数据本身**，只做以下校验：

\begin{itemize}
  \item YAML 语法正确
  \item 必填字段齐全（schema 校验）
  \item \yamlkey{project.years} 范围合法（start <= end）
  \item Reference 数据集名能在 registry 里找到（含 \texttt{\_LowRes} / \texttt{\_MidRes} 后缀自动解析）
  \item 解析过程中标注 provenance（数据集元信息的可信度）
\end{itemize}

成功输出：

\begin{exampleBox}[check 通过的输出]
\begin{minted}{text}
Config validation:
  ✓ YAML syntax valid
  ✓ Schema validation passed
  ✓ Year range [2010, 2014] valid

Reference data (2 sources):
  ✓ Evapotranspiration → GLEAM_v4.2a_LowRes (grid, Month, 0.5°)
  ✓ Total_Runoff → GRUN_ENSEMBLE_LowRes (grid, Month, 0.5°)

Simulation data (1 models):
  ✓ CoLM2024 (model: CoLM2024, root: /data/colm2024)

Options:
  Time alignment: intersection
  Unified mask: True
  Comparison: True
  Statistics: False

✓ Config valid (2 variables, 2 references, 1 simulations). Ready to run.
\end{minted}
\end{exampleBox}

\subsection{常见 check 失败}

\textbf{Reference 数据集找不到}：

\begin{minted}{text}
  ✗ Evapotranspiration → GLEAM_v999: not in registry
\end{minted}

修法：用 \cli{openbench data list} 查看可用名称，或检查 \texttt{\_LowRes}/\texttt{\_MidRes} 后缀。

\textbf{Provenance 警告}：

\begin{minted}{text}
  ✓ Evapotranspiration → GLEAM_v4.2a_LowRes
    ⚠ tim_res: Month (default — not confirmed from NC or profile)
\end{minted}

意思是注册表里数据集的"时间分辨率"字段是默认假设值，没在元数据里被验证。如果你启用 \yamlkey{project.strict_reference: true}，这种警告会变成错误。详见第 6 章。

\textbf{歧义解析}：

\begin{minted}{text}
  ✗ Soil_Moisture → ESA_CCI_SM
    Multiple matches: ESA_CCI_SM_v9.0_LowRes, ESA_CCI_SM_v9.0_MidRes
\end{minted}

修法：把 \yamlkey{reference.sources} 里的简称改为完整的 \texttt{\_LowRes} 或 \texttt{\_MidRes} 全名。

\section{步骤 3：运行评估（\cli{openbench run}）}

\begin{minted}{bash}
openbench run openbench.yaml
\end{minted}

这是真正读数据、做计算、出图的命令。典型输出：

\begin{exampleBox}[run 成功的输出]
\begin{minted}{text}
Running evaluation: demo
  Phase: preprocess
    Loading reference: GLEAM_v4.2a_LowRes
    Regridding to 0.5° ...
    ...
  Phase: evaluation
    Computing metrics for Evapotranspiration ...
    Computing metrics for Total_Runoff ...
  Phase: comparison
    Cross-model comparison enabled, 1 model(s)
  Phase: visualization
    Writing 23 figures to output/demo/figures/
  Phase: report
    Generated output/demo/report.html

✓ Evaluation complete
  Output: output/demo
  Variables: 2
  Simulations: 1
\end{minted}
\end{exampleBox}

\subsection{常用运行选项}

\begin{minted}{bash}
# 先做一次 dry-run（验证 + 不实际计算）
openbench run openbench.yaml --dry-run

# 限制只跑某几个变量（覆盖配置文件）
openbench run openbench.yaml --variables Evapotranspiration

# 覆盖并行核数（覆盖配置文件 project.num_cores）
openbench run openbench.yaml --cores 8

# 只重做 comparison 阶段，跳过评估（前提：评估输出已存在）
openbench run openbench.yaml --comparison-only

# 输出中间 runner / namelist 配置到 output/<run>/debug/
openbench run openbench.yaml --dump-config
\end{minted}

\warninline{\cli{--remote} 标志当前在 CLI 中尚未实现，传入会立刻退出并提示用 GUI。运维卷第 3 章会详细介绍远程执行的现状与替代方案。}

\section{步骤 4：看输出目录}

评估完成后，\file{output/demo/} 大致结构：

\begin{minted}{text}
output/demo/
├── config.yaml              # 本次评估使用的完整配置（含解析后默认值）
├── debug/                   # 仅当 --dump-config 时存在
│   ├── runner_config.yaml
│   ├── main_nl.yaml
│   ├── ref_nml.yaml
│   ├── sim_nml.yaml
│   └── fig_nml.yaml
├── data/                    # 预处理后的中间数据（zarr 缓存）
├── evaluation/
│   ├── Evapotranspiration/
│   │   ├── CoLM2024_metrics.csv
│   │   ├── CoLM2024_scores.csv
│   │   └── ...
│   └── Total_Runoff/...
├── comparison/              # 多模型对比表（>= 2 个 simulation 时）
├── figures/                 # 全部可视化产出
│   ├── geo/                 # 地理分布图
│   ├── portrait/            # portrait plot
│   ├── radar/               # 雷达图
│   └── ...
└── report.html              # 总结报告（[report] extra 启用时）
\end{minted}

第 7 章会逐项解释这些目录的内容。

\section{下一步}

\begin{itemize}
  \item 想了解配置文件的全部字段：去第 3 章
  \item 想看 56 个评估变量与 67 个 reference 数据集的全貌：去第 4 章
  \item 想知道 \openbench{} 在每个阶段做什么：去第 6 章
  \item 想用 GUI 而非命令行：去第 8 章
\end{itemize}
